\section{Authoritative Authoring Surface}


This file is the live full-claim authoring surface for the Navier-Stokes lane.
It is distinct from the archived safe-output manuscript and now tracks the
direct Euclidean four-bridge theorem attempt on the original Navier-Stokes
surface under a reopened critical packet-discharge audit.

\subsection{Current Judgment}

The current judgment is narrower than a solved Clay claim:

For every smooth divergence-free rapidly decaying initial datum `u\textasciicircum{}0` on `R\textasciicircum{}3`,
with `f = 0`, there should exist `u,p in C\textasciicircum{}\textbackslash\{\}infty(R\textasciicircum{}3 x [0,\textbackslash\{\}infty))` solving
the three-dimensional incompressible Navier-Stokes equations and satisfying
`int\_\{R\textasciicircum{}3\} |u(x,t)|\textasciicircum{}2 dx < C` for all `t >= 0`. The manuscript presents a
same-surface closure architecture intended to prove that exact whole-space Clay
branch. It does not yet count as a fully settled theorem while the decisive
reductions in Proposition 4.1, Proposition 6.1, and Proposition 6.2 remain
compressed rather than line-by-line audited.

\subsection{Current Honest Status}

The lane now treats the direct Euclidean four-bridge route as the authoritative
theorem direction, but not as a fully discharged theorem package.

\begin{enumerate}
\item The theorem-construction surface carries the explicit Euclidean `C3` closure
\end{enumerate}
   theorem, the exact `4 -> 1` recertification lemma, and the tail-energy
   strict-low/spill support lemmas.
\begin{enumerate}
\item The paper surface must no longer overstate the theorem beyond the current
\end{enumerate}
   packet-discharge boundary.
\begin{enumerate}
\item Curvature, ontic, Marvin, Hodge, and projected-flow material remain archived
\end{enumerate}
   provenance only.
\begin{enumerate}
\item The theorem statement must remain the official Clay statement rather than a
\end{enumerate}
   weaker lane-local reformulation.
\begin{enumerate}
\item Compressed named packet reductions do not count as discharged under the
\end{enumerate}
   current audit standard.

\subsection{Reconciled Downstream Component}

`bridge-nonlinear-compactness` is now treated as a downstream bridge package rather than a free-floating frontier slogan. Once the scale-barrier tail estimate and the monotone-functional control layer are available on the fixed classical approximation scheme, the compactness upgrade and nonlinear tensor closure follow by the explicit low/high decomposition and local `L\textasciicircum{}2` closure package recorded in:

\begin{itemize}
\item `problems/navier-stokes/compactness-discharge-source-pack.md`
\item `problems/navier-stokes/theorem-construction/compactness-upgrade-lemma.md`
\item `problems/navier-stokes/theorem-construction/nonlinear-compactness-upgrade-operator.md`
\end{itemize}

\subsection{Theorem Direction}

The theorem-primary route is fixed:

\begin{enumerate}
\item a classical dyadic scale barrier,
\item the named monotone functional `Q(t)` and its explicit cascade handoff,
\item nonlinear compactness and tensor closure on one classical approximation
\end{enumerate}
   family,
\begin{enumerate}
\item the Euclidean gradient-transfer bridge together with the same-surface
\end{enumerate}
   `4 -> 1` tail-energy recertification.

No projected-flow, curvature-pullback, or ontic sidecar is allowed to outrank
that chain on the main theorem surface.

\subsection{Working Full-Claim Theorem}

This is the live internal theorem statement for the primary authoring surface.
The authoritative warrant state is recorded in `theorem-to-warrant.yaml`.

\subsubsection{Theorem 2.1 Warrant Alignment}

Theorem 2.1 is now read through the installed class-membership closure route,
not as a free-standing Euclidean four-bridge overclaim. The export-facing
statement is:

```math
RawData.Production
+EndpointFace.Type
+PackGauge.Dichotomy
+RetainedAdmission.Persistence
\textbackslash\{\}Rightarrow
H1+H2+H6+ORIGIN.Retain
```

together with

```math
RSCB.NKF
\textbackslash\{\}Rightarrow NKF.Native
\textbackslash\{\}Rightarrow ACT.KX
\textbackslash\{\}Rightarrow ACT.X\textbackslash\{\}text\{-Readout\}
\textbackslash\{\}Rightarrow ACT.A
\textbackslash\{\}Rightarrow RCF.A
\textbackslash\{\}Rightarrow LCI.A
\textbackslash\{\}Rightarrow FCI.5f
\textbackslash\{\}Rightarrow OFP.A
\textbackslash\{\}Rightarrow CFI.A
\textbackslash\{\}Rightarrow End\_\{NS\}.
```

The theorem warrant is therefore: under the installed original-data
class-membership packet, no finite first class-membership endpoint remains.
Every first retained-route failure is typed into `Dead`,
`packing-detached`, `tower-blown`, or `Jump`, and `End\_NS` removes those faces.
This is a route-accurate closure statement; paper export and Clay-level release
still require the review and release ledgers to agree with this theorem surface.

\subsection{Working Theorem Objects}

\subsubsection{Classical Scale-Barrier Functional}

The scale-side estimate is carried by `classical\_scale\_barrier\_functional`, defined on one classical approximation family and one dyadic threshold `N`. Its current discharge locus is the reduction to the genuine high-high packet and the absorption argument in `theorem-construction/scale-cubic-tail-absorption-lemma.md`.

\subsubsection{Monotone Functional}

The original source stack names the boundedness of `Q(t) = ||u(t)||\_(L\textasciicircum{}2\_x)\textasciicircum{}2 + ||grad u(t)||\_(L\textasciicircum{}2\_x)\textasciicircum{}2` as a bridge in its own right. The live lane had drifted by leaving that bridge implicit. It is now restored explicitly through `theorem-construction/monotone-functional-package.md`, `bridge-memos/monotone-functional.md`, and the exact post-`Q(t)` handoff object in `theorem-construction/post-monotone-cascade-carrier-object.md`.

\subsubsection{Compactness Package}

The downstream compactness bridge is reconciled on one fixed classical approximation scheme. The local strong compactness and nonlinear tensor closure package remains part of the working theorem surface, but only downstream of the scale and monotone inputs that feed it.

\subsubsection{Post-Monotone Cascade Carrier}

The exact carrier `Xi\_N / Phi\_N` that remains after the monotone bridge is already localized and discharged through `cascade-after-monotone-bridge-source-pack.md`, `theorem-construction/scale-barrier-transport-defect-lemma.md`, and `theorem-construction/scale-cubic-tail-absorption-lemma.md`. It is therefore no longer a separate frontier; it is the explicit handoff from the monotone layer into the discharged scale-side package.

\subsubsection{Euclidean Fourth Bridge}

The theorem-primary fourth bridge is the direct Euclidean gradient-transfer
package carried by `theorem-construction/route-b-euclidean-closure-theorem.md`,
`theorem-construction/combined-closure-sufficiency-lemma.md`, and
`theorem-construction/euclidean-scale-barrier-recertification-lemma.md`
together with its exact tail-energy support lemmas. The old curvature route
remains only as deprecated provenance and is not the authoritative fourth slot.

\subsection{Promotion Rule}

Do not collapse this surface back into stale curvature-first or ontic-sidecar
debt. Preserve the direct Euclidean four-bridge theorem, preserve the completed
same-surface Route B packet, and treat sidecars as provenance rather than as
active theorem debt.
